\section{A Complete Application with Jolie}
\label{sec:complete}

In this section, we will describe a complete application written in Jolie. The
application will contain several services and will be written using best
practices from before the module and package system.

Along the way, several patterns used for re-using code will be highlighted.
This will serve as the core motivation for our solution. The observations made
in this section are drawn from the Jolie Documentation\autocite{JOLB}, the
Leonardo\footnote{A web server written in Jolie by the Jolie team. Used to
    power the official Jolie website} web server\autocite{LEO}, and the
    source-code for the Jolie website\autocite{LEOB}.

\subsection{Architecture}

% Things we wish to cover in this section:
%
%   - Includes and why they will be problematic
%   - Configuring location of ports, both external and embedding of stuff
%   - Project structure, and why this is a problem
%   - Parameter configuration

The application we'll be building is a very simple calculator system. The
architecture of our application is shown in Figure \ref{fig:complete_arch}.

The calculator service will provide various operations for applying an operator
on a sequence of numbers. In order to perform this, the action of applying
these operators, the calculator will contact other microservices. This example
is obviously contrived. Writing services for such simplistic services in the
real world would be considered an anti-pattern. However, this example allows us
to see services interact with each other without thinking that much about the
business logic of the services, which in this case is fairly irrelevant.

\begin{figure}[H]
\centering
\includegraphics[width=1.0\textwidth]{background/calc_system.eps}
\caption{The Architecture of a Simple Microservice System}
\label{fig:complete_arch}
\end{figure}

\subsection{Implementing the Calculator Service}

We will keep our focus on the \txtl{calculator} service, and its interactions
with other services. Putting together the stuff learned from the previous
sections, we can quickly create an input port for the service with the
appropriate interface.

It is considered best practice to place the public interfaces that a service
exposes in its own separate file. Files intended for other services to include
typically have the file extension \txtl{.iol} as opposed to \txtl{.ol}. There
is no technical difference between the two, but it allows for the developer to
more easily express intent. Files that are included can be placed in the
\txtl{include} directory. This directory is always implicitly added to the
search path.

Thus in order to implement our calculator service we create two files. One file
for the service implementation (\verb!calculator.ol!), and another which can be
used by other services (\verb!calculator.iol!). These files are shown in
Listing \ref{lst:simple_start}.

\begin{listing}[H]
\begin{multicols}{2}

\begin{minted}[breaklines,fontsize=\footnotesize]{jolie}
// calculator.ol
include "console.iol"

inputPort Calculator {
    Location: "socket://localhost:12345"
    Protocol: sodep
    Interfaces: ICalculator
}

main {
    [sum(request)(response) {
        println@Console("Implementation goes here")()
    }]

    [product(request)(response) {
        println@Console("Implementation goes here")()
    }]
}
\end{minted}

\columnbreak

\begin{minted}[breaklines,fontsize=\footnotesize]{jolie}
// calculator.iol
type Numbers: void {
    .numbers[2, *]: int
}

interface ICalculator {
    RequestResponse:
        sum(Numbers)(int),
        product(Numbers)(int)
}
\end{minted}
\end{multicols}
\caption{Interface and implementation of the \txtl{calculator} service}
\label{lst:simple_start}
\end{listing}

\begin{observation}

The interfaces and types of a service are typically separated into their own
files. These files typically have the extension \verb!.iol! to indicate that it
does not contain a service implementation, but is rather intended for inclusion
by a service which requires it.

These files are typically put in the \verb!include! directory. This
directory is always added to the search path natively by the Jolie engine.

\end{observation}

The operations which the calculator exposes needs to collaborate with the other
services. In order for us to speak to them, they need output ports.  The output
ports need interfaces, just like the calculator service these should be exposed
from a dedicated file.

Secondly, these output port needs to be reached. The port can either be
embedded or we can provide external bindings to it. For output ports, we may
change this binding dynamically at runtime. This is unlike input ports which
must be ready at deployment time.

Binding an output port to an external service, is done by setting the
\joliel{Location} and \joliel{Protocol} properties. For example, to make
\joliel{Multiplication} bind to an external service we could define it as shown
in Listing \ref{lst:mult_ext_port}.

\begin{listing}[H]
\begin{minted}{jolie}
outputPort Multiplication {
    Location: "socket://mult.example.com:443"
    Protocol: https
}
\end{minted}
\caption{The \joliel{Multiplication} port is bound to an external service}
\label{lst:mult_ext_port}
\end{listing}

It is possible to define embeddings of services in the \joliel{embedded} block.
For our desired deployment, we wanted to embed the addition service inside of
the calculator service. Assuming that the addition service is written as a
Jolie service, and its service is implemented in \txtl{addition.ol}, then we
may create an embedding as shown in Listing \ref{lst:simple_embedding}. Just
like in the case of the external services, the input port of the receiving
service \emph{must match}. In the case of embedded service, there must be an
input port listening on the \joliel{"local"} location.

\begin{listing}[H]
\begin{minted}{jolie}
embedded {
    Jolie:
        "addition.ol" in Addition
}
\end{minted}

\caption{Embedding the \mintinline{jolie}{addition} service in the
    \mintinline{jolie}{Addition} output port}

\label{lst:simple_embedding}

\end{listing}

Quite often the location of an input port is considered a deployment problem.
We see this quite clearly in the case of embedding a service. All solutions to
change the deployment of a service in Jolie, all require us to \emph{modify the
    source code of a service}.  Best practices in Jolie attempt to make this
    less of a problem by including a configuration file which contains
    constants. The \txtl{addition} service, might include a file called
    \verb!addition.iol! which sets up the location and protocol of the service.
    An example of this is shown in Listing \ref{lst:include_as_conf}.

\begin{listing}[H]
\begin{multicols}{2}

\begin{minted}[breaklines,fontsize=\footnotesize]{jolie}
// addition.ol
include "addition.iol"

inputPort Addition {
    Location: ADDIITON_LOCATION
    Protocol: ADDIITON_PROTOCOL
}
\end{minted}

\columnbreak

\begin{minted}[breaklines,fontsize=\footnotesize]{jolie}
// addition.iol
constants {
    ADDIITON_LOCATION = "local"
    ADDIITON_PROTOCOL = sodep
}
\end{minted}

\end{multicols}

\caption{A common Jolie practice for solving configuration of a service, is to
    include a file containing constants with the desired configuration.}

\label{lst:include_as_conf}

\end{listing}

\begin{observation}

Most configuration of Jolie services is done via constants. Jolie constants,
different from most other languages, can be simple literal types, or they
may even contain identifiers. This makes them rather versatile in what
they can configure.

The constants are typically placed in a separate source file, which is
included by the service requiring the configuration.

\end{observation}

In Listing \ref{lst:simple_start} the output port \txtl{Console} is used. This
output port isn't visibly defined in the source code we see, it is instead
defined in the included file \txtl{console.iol}. This is common practice in
Jolie. This is especially used for services that work in a library-like
fashion, that is not intended as a stand-alone service. This pattern is used
for almost all services in the Jolie standard library\autocite{JOLB}.

\begin{observation}

Services intended to be used as libraries are often contained in a single
\verb!.iol! file. This file contains everything required to set it up, included
interfaces, types and an embedded output port.

This makes it very easy to use the service, however, it also makes it
impossible to bind to this service externally, without first getting an
embedding. This isn't possible since we cannot include the interfaces and
types by themselves.

\end{observation}

With the output ports correctly configured, we may now implement the actual
business logic for our calculator. This might look like shown in Listing
\ref{lst:op_impl}.

\begin{listing}[H]
\begin{minted}{jolie}
[sum(request)(response) {
    total = 0;
    for (i = 0, i < #request.numbers, i++) {
        // Add current number with total, and store result in total
        add@Addition({
            .a = total,
            .b =  request.numbers[i]
        })(total)
    };
    response = total
}]
\end{minted}
\caption{Implementing the checkout operation}
\label{lst:op_impl}
\end{listing}

Finally, we want to reflect slightly on the file structure that this service
has. In order to use external services, we had to include files which contain
these interfaces. This version is worse when an embedding is desired, since the
entire service with its implementation is suddenly required.

The files from these external services are also hard to manage. It isn't
possible to simply move the source code of a service into its own directory.
This is not possible since includes, unlike in most other languages, are not
relative to the file performing the include, but rather relative to the project
level root. Thus if we wish to move the service to a new directory, all
includes in the source code would have to be updated. A common result of this
is that every single file gets dumped into the project level root. When all
files are stored in the same directory, we also get a much larger chance of
having name conflicts. As a result, names tend to become rather large, to make
it less likely that a conflict occurs.

The final file structure of the calculator service is shown in Listing
\ref{lst:file_structure}.

\begin{listing}[H]
\begin{minted}{text}
.
|-- addition.iol
|-- addition.ol
|-- calculator.ol
|-- include
|   `-- calculator.iol
`-- multiplication.iol
\end{minted}
\caption{File structure of the calculator service}
\label{lst:file_structure}
\end{listing}

\begin{observation}

    In order to collaborate with other services, it is often needed to manually
    copy files into the service. Which files are required depend on the
    deployment. When an embedding is required, we will need the entire service
    (and any services it may depend on itself). If we just need to interface
    with it, we can simply include the files required for its interface files.

    Structuring the files is problematic, due to \verb!include! statements not
    being relative to file performing inclusion.

\end{observation}

